\begin{theorem}[unitary-slice 边界审计的三轴不可约分解]\label{thm:conclusion-unitary-slice-boundary-three-axis-irreducibility}
设
\[
\mathfrak M_{\partial},
\qquad
\mathfrak Z_{\partial},
\qquad
\mathfrak E_{\partial}
\]
分别表示 unitary-slice 可见域上的边界模长通道、盘内零点通道与端点载荷通道。若某一边界审计函子 \(\mathfrak B\) 对这三类数据同时 faithful，则 \(\mathfrak B\) 必经由积对象
\[
\mathfrak M_{\partial}\times \mathfrak Z_{\partial}\times \mathfrak E_{\partial}
\]
因子化；并且该分解是最小的：去掉任一因子后，所得审计都不再完备。
\end{theorem}
\leanverified{paper\_conclusion\_unitary\_slice\_boundary\_three\_axis\_irreducibility}
\begin{proof}
正文关于 unitary-slice 正交切片的结果已经把边界层可见信息分裂为三条彼此不可替代的通道：可见模长由 pure-phase 接口承载，盘内零点由内因子与 Jensen 缺陷几何承载，而端点载荷则由端点热核探针独立读出。特别地，定理 \ref{thm:conclusion-visible-channel-anomaly-finite-presentation} 与定理 \ref{thm:conclusion-toeplitz-endpoint-atom-dyadic-joint-verification} 已表明：盘内零点与端点原子分别拥有各自的有限否证器与独立抽取接口，二者都不能由单纯的边界模长证书替代。

因此，若删去 \(\mathfrak Z_{\partial}\)，则盘内零点事件无法恢复；若删去 \(\mathfrak E_{\partial}\)，则端点载荷与端点原子质量不可恢复；若删去 \(\mathfrak M_{\partial}\)，则 pure-phase 边界模长审计本身退化。于是任何对三类数据同时 faithful 的边界审计都只能经由上述三轴积对象因子化，且不存在二轴压缩后的完备实现。证毕。
\end{proof}
